\documentclass[10pt,a4]{examenEPN0}
\usepackage{comandosEPN}

\usepackage{empheq}

\newcommand{\itemcp}{\stepcounter{enumi}\item[ \textcolor{red}{\theenumi.$^*$}] }

\newcommand{\iitemcp}{\stepcounter{enumii}\item[ \textcolor{red}{\theenumii)$^*$}] }


\materia{Álgebra Lineal}
%\autor{Departamento de Formación Básica}
\autor{GR \hspace{1cm}.}
%\fecha{Semestre 2026-A}
\fecha{Nombre: \hspace{9cm} Cédula:}
\logouno[1.6cm]{Logo_EPN.eps}
\logodos[2.6cm]{Logo_DFB.eps}

\newcommand{\la}{\lambda}

%Datos
\newcommand{\datosAlgebra}[2]{
\EnExamen
{
\vspace{5mm}
    \noindent
    Nombres: \hrulefill\
    Carrera: \rule{4.5cm}{0.4pt}\
    \ifthenelse{\equal{#1}{}}%
    {Paralelo: \underline{\parbox{1.8cm}{\parbox{1cm}{}}}\\[3mm]}
    {Paralelo: \underline{\parbox{1.8cm}{\centering #1}}\\[3mm]}
    C\'edula: \rule{5cm}{0.4pt}\
    Firma: \hrulefill\
    Fecha: \underline{\parbox{2.1cm}{\centering #2}}
}}


% \newcommand{\indicacionesPrueba}{%
% \EnExamen{
% \begin{indicaciones}
%     \item La prueba tiene una duraci\'on  de \textbf{120 minutos.}
%     \item No se permite la utilizaci\'on de dispositivos electr\'onicos (celular, tablets, calculadora, etc.)
%     \item Leer detenidamente los enunciados, \textbf{justificar la respuesta para acceder al puntaje total}.
%     \item La prueba es completamente individual. Los estudiantes deberán dejar sus teléfonos celulares guardados y no podrán tenerlos consigo (por ejemplo, en los bolsillos de su ropa) durante toda la evaluación. En caso de detectarse algún acto de deshonestidad académica o el incumplimiento de esta norma, no podrán seguir rindiendo la prueba y su calificación será \textbf{cero}.
% \end{indicaciones}
% \pagestyle{empty}
% }}

\newcommand{\indicacionesPrueba}{%
\EnExamen{
\begin{indicaciones}
    \item La prueba tiene una duraci\'on  de \textbf{60 minutos.}
    \item Leer detenidamente los enunciados, \textbf{justificar la respuesta} para acceder al puntaje total.
    \item No se permite la utilizaci\'on de dispositivos electr\'onicos (celular, tablets, calculadora, etc.)
    \item La prueba es completamente individual
    \item No se permite portar teléfonos celulares, relojes inteligentes ni otros dispositivos electrónicos durante el examen, aun cuando permanezcan en los bolsillos. Estos dispositivos deberán guardarse en la mochila antes del inicio de la evaluación y esta deberá mantenerse cerrada durante todo el examen.
    \item He leído las indicaciones y acepto que el incumplimiento de estas se penaliza con \textbf{cero}. Firma:...................................

    
\end{indicaciones}
\pagestyle{empty}
}}

\newcommand{\indicacionesExamen}{
\EnExamen{\begin{indicaciones}
    \item La evaluación tiene una duraci\'on  de \textbf{120 minutos.}
    \item No se permite la utilizaci\'on de dispositivos electr\'onicos (celular, tablets, calculadora, etc.)
    %\item Leer detenidamente los enunciados, justificar la respuesta para acceder al puntaje total.
    \item El examen es completamente individual. En caso de detectarse alg\'un acto de deshonestidad acad\'emica no podr\'a seguir rindiendo la prueba y su calificaci\'on ser\'a \textbf{cero}.
\end{indicaciones}
}\pagestyle{empty}
}

%% Imprimir las una hoja A4, dos veces en A5
\newcommand{\imprimirDoble}
{%
\usepackage{pgfmorepages}
\pgfmorepagesloadextralayouts
\pgfpagesuselayout{repeated 2-up}[a4paper,landscape]
\thispagestyle{empty}
}


%Este sí se puede comentar
\newenvironment{solucion}{\color{blue}\proof[Solución]}{\endproof\color{black}}




% \renewcommand{\@autor}{
% ---\\
% \phantom{.} 
% ...}
% \makeatother

\examen{Prueba 02: solucionario}
\makeatletter
\renewcommand{\@autor}{Departamento de Formación Básica}
\renewcommand{\@fecha}{GR5M -- Miércoles de 7 a 9 -- 04/11/2026}
\makeatother

% Soluciones de Miercoles_7_9_Prueba02.tex (20 puntos).
\begin{document}
\encabezado

\noindent\textbf{Criterio general de calificación.} Se asigna crédito parcial
por los pasos correctos del procedimiento. Un error aritmético se penaliza
en el paso donde se origina; los pasos posteriores se valoran por su
coherencia con los resultados previos, sin repetir la misma penalización.
Se aceptan procedimientos equivalentes que respeten el método solicitado.

\section*{1. Mezcla de fertilizantes (6 puntos)}
\noindent\textbf{Rúbrica (6 puntos).}
\begin{itemize}[itemsep=2pt]
    \item Definir las tres variables con sus unidades: 0,5 puntos.
    \item Plantear las ecuaciones de los nutrientes: 1,5 puntos
          (0,5 por cada ecuación).
    \item Identificar correctamente $A$, $x$ y $b$: 1 punto.
    \item Construir $[A\mid I_3]$: 0,25 puntos.
    \item Escalonar con operaciones por filas identificadas y correctas: 1 punto.
    \item Reducir hasta $[I_3\mid A^{-1}]$, identificando las operaciones:
          1 punto.
    \item Calcular $x=A^{-1}b$: 0,75 puntos.
\end{itemize}

\noindent\textbf{a) Variables y modelo matricial (3 puntos).}
Sean $x_A$, $x_B$ y $x_C$ las cantidades de sacos de los abonos A, B y C,
respectivamente. Por el contexto, las tres cantidades deben ser no negativas.
Las ecuaciones para los nutrientes, medidos en kilogramos, son
\[
    \left\{\begin{aligned}
        x_B+2x_C&=7 &&\text{(nitrógeno)},\\
        2x_A+x_B+3x_C&=14 &&\text{(fósforo)},\\
        x_A+x_B+x_C&=6 &&\text{(potasio)}.
    \end{aligned}\right.
\]
La representación $Ax=b$ queda determinada por
\[
    A=\begin{pmatrix}0&1&2\\2&1&3\\1&1&1\end{pmatrix},\qquad
    x=\begin{pmatrix}x_A\\x_B\\x_C\end{pmatrix},\qquad
    b=\begin{pmatrix}7\\14\\6\end{pmatrix}.
\]

\noindent\textbf{b) Método de la matriz inversa (3 puntos).}
Para calcular la inversa mediante Gauss--Jordan, colocamos la identidad
a la derecha de $A$ y aplicamos cada operación a toda la fila de la
matriz ampliada $[A\mid I_3]$.

\noindent\textit{Escalonamiento.} Primero intercambiamos las filas 1 y 3
para obtener un pivote igual a 1:
\[
    \left[\begin{array}{ccc|ccc}
        0&1&2&1&0&0\\2&1&3&0&1&0\\1&1&1&0&0&1
    \end{array}\right]
    \xrightarrow{F_1\leftrightarrow F_3}
    \left[\begin{array}{ccc|ccc}
        1&1&1&0&0&1\\2&1&3&0&1&0\\0&1&2&1&0&0
    \end{array}\right].
\]
Anulamos la entrada debajo del primer pivote y normalizamos la segunda fila:
\[
    \left[\begin{array}{ccc|ccc}
        1&1&1&0&0&1\\2&1&3&0&1&0\\0&1&2&1&0&0
    \end{array}\right]
    \xrightarrow{F_2\leftarrow F_2-2F_1}
    \left[\begin{array}{ccc|ccc}
        1&1&1&0&0&1\\0&-1&1&0&1&-2\\0&1&2&1&0&0
    \end{array}\right]
    \xrightarrow{F_2\leftarrow -F_2}
    \left[\begin{array}{ccc|ccc}
        1&1&1&0&0&1\\0&1&-1&0&-1&2\\0&1&2&1&0&0
    \end{array}\right].
\]
Ahora anulamos la entrada debajo del segundo pivote:
\[
    \left[\begin{array}{ccc|ccc}
        1&1&1&0&0&1\\0&1&-1&0&-1&2\\0&1&2&1&0&0
    \end{array}\right]
    \xrightarrow{F_3\leftarrow F_3-F_2}
    \left[\begin{array}{ccc|ccc}
        1&1&1&0&0&1\\0&1&-1&0&-1&2\\0&0&3&1&1&-2
    \end{array}\right].
\]
El bloque izquierdo está en forma escalonada y tiene tres pivotes.

\noindent\textit{Reducción.} Normalizamos el tercer pivote y anulamos
la entrada situada sobre él en la segunda fila:
\[
    \left[\begin{array}{ccc|ccc}
        1&1&1&0&0&1\\0&1&-1&0&-1&2\\0&0&3&1&1&-2
    \end{array}\right]
    \xrightarrow{F_3\leftarrow \frac13 F_3}
    \left[\begin{array}{ccc|ccc}
        1&1&1&0&0&1\\0&1&-1&0&-1&2\\0&0&1&\frac13&\frac13&-\frac23
    \end{array}\right]
    \xrightarrow{F_2\leftarrow F_2+F_3}
    \left[\begin{array}{ccc|ccc}
        1&1&1&0&0&1\\0&1&0&\frac13&-\frac23&\frac43\\0&0&1&\frac13&\frac13&-\frac23
    \end{array}\right].
\]
Finalmente, anulamos las entradas sobre los pivotes en la primera fila:
\[
    \left[\begin{array}{ccc|ccc}
        1&1&1&0&0&1\\0&1&0&\frac13&-\frac23&\frac43\\0&0&1&\frac13&\frac13&-\frac23
    \end{array}\right]
    \xrightarrow{F_1\leftarrow F_1-F_3}
    \left[\begin{array}{ccc|ccc}
        1&1&0&-\frac13&-\frac13&\frac53\\0&1&0&\frac13&-\frac23&\frac43\\0&0&1&\frac13&\frac13&-\frac23
    \end{array}\right]
    \xrightarrow{F_1\leftarrow F_1-F_2}
    \left[\begin{array}{ccc|ccc}
        1&0&0&-\frac23&\frac13&\frac13\\0&1&0&\frac13&-\frac23&\frac43\\0&0&1&\frac13&\frac13&-\frac23
    \end{array}\right].
\]
Hemos obtenido $[I_3\mid A^{-1}]$, por lo que $A$ es invertible y
\[
    A^{-1}=\frac{1}{3}\begin{pmatrix}-2&1&1\\1&-2&4\\1&1&-2\end{pmatrix},
    \qquad
    x=A^{-1}b=\frac{1}{3}
    \begin{pmatrix}-14+14+6\\7-28+24\\7+14-12\end{pmatrix}
    =\begin{pmatrix}2\\1\\3\end{pmatrix}.
\]
\textbf{Se necesitan 2 sacos de A, 1 saco de B y 3 sacos de C.}
Las cantidades son no negativas y verifican las exigencias:
\[
    1+2(3)=7,\qquad 2(2)+1+3(3)=14,\qquad 2+1+3=6.
\]

\section*{2. Construcción de una matriz diagonal (6 puntos)}
\noindent\textbf{Rúbrica (6 puntos).}
\begin{itemize}[itemsep=2pt]
    \item Obtener las tres componentes de $x$ mediante Cramer: 1,5 puntos
          (0,5 por cada componente).
    \item Proponer una matriz diagonal con determinante 3: 1 punto.
    \item Calcular $b=Ax$: 1,5 puntos (0,5 por cada componente).
    \item Verificar los cuatro determinantes solicitados: 2 puntos
          (0,5 por cada determinante).
\end{itemize}

Denotamos por $A^{(i)}$ la matriz obtenida al reemplazar la columna $i$
de $A$ por $b$. Como $\det(A)=3\ne0$, la regla de Cramer da
\[
    x_1=\frac{-9}{3}=-3,\qquad
    x_2=\frac{12}{3}=4,\qquad
    x_3=\frac{6}{3}=2.
\]
Elegimos una matriz diagonal cuyo producto de entradas diagonales sea 3:
\[
    \boxed{A=\begin{pmatrix}1&0&0\\0&1&0\\0&0&3\end{pmatrix},
    \qquad b=A\begin{pmatrix}-3\\4\\2\end{pmatrix}
    =\begin{pmatrix}-3\\4\\6\end{pmatrix}.}
\]
Verificamos los cuatro determinantes solicitados:
\[
    \det(A)=1\cdot1\cdot3=3,
    \qquad
    \det\begin{pmatrix}-3&0&0\\4&1&0\\6&0&3\end{pmatrix}=-9,
\]
\[
    \det\begin{pmatrix}1&-3&0\\0&4&0\\0&6&3\end{pmatrix}=12,
    \qquad
    \det\begin{pmatrix}1&0&-3\\0&1&4\\0&0&6\end{pmatrix}=6.
\]
Este ejemplo satisface todas las condiciones; la elección de $A$ y $b$
no es única.

\newpage
\section*{3. Resolución mediante la regla de Cramer (6 puntos)}
\noindent\textbf{Rúbrica (6 puntos).}
\begin{itemize}[itemsep=2pt]
    \item Identificar $A$ y $b$: 0,5 puntos.
    \item Calcular $\det(A)$ y justificar que Cramer es aplicable: 1 punto.
    \item Construir $A^{(1)}$, $A^{(2)}$ y $A^{(3)}$: 1,5 puntos
          (0,5 por cada matriz).
    \item Calcular $\det(A^{(1)})$, $\det(A^{(2)})$ y $\det(A^{(3)})$:
          1,5 puntos (0,5 por cada determinante).
    \item Aplicar los tres cocientes de Cramer y obtener $(x,y,z)$:
          1,5 puntos (0,5 por cada incógnita).
\end{itemize}

El sistema dado tiene matriz de coeficientes y término independiente
\[
    A=\begin{pmatrix}1&2&1\\2&-1&1\\3&1&-2\end{pmatrix},
    \qquad b=\begin{pmatrix}4\\2\\2\end{pmatrix}.
\]
Desarrollamos los determinantes por la primera fila:
\[
    \det(A)
    =1(2-1)-2(-4-3)+1(2+3)=20.
\]
Como $\det(A)\ne0$, el sistema tiene solución única. Denotamos por
$A^{(i)}$ la matriz obtenida al reemplazar la columna $i$ de $A$ por $b$:
\[
    A^{(1)}=\begin{pmatrix}4&2&1\\2&-1&1\\2&1&-2\end{pmatrix},\qquad
    A^{(2)}=\begin{pmatrix}1&4&1\\2&2&1\\3&2&-2\end{pmatrix},\qquad
    A^{(3)}=\begin{pmatrix}1&2&4\\2&-1&2\\3&1&2\end{pmatrix}.
\]
Sus determinantes son
\begin{align*}
    \det(A^{(1)})&=\det\begin{pmatrix}4&2&1\\2&-1&1\\2&1&-2\end{pmatrix}
        =4(2-1)-2(-4-2)+(2+2)=20,\\[2mm]
    \det(A^{(2)})&=\det\begin{pmatrix}1&4&1\\2&2&1\\3&2&-2\end{pmatrix}
        =(-4-2)-4(-4-3)+(4-6)=20,\\[2mm]
    \det(A^{(3)})&=\det\begin{pmatrix}1&2&4\\2&-1&2\\3&1&2\end{pmatrix}
        =(-2-2)-2(4-6)+4(2+3)=20.
\end{align*}
La regla de Cramer implica
\[
    x=\frac{\det(A^{(1)})}{\det(A)}=\frac{20}{20}=1,\qquad
    y=\frac{\det(A^{(2)})}{\det(A)}=\frac{20}{20}=1,\qquad
    z=\frac{\det(A^{(3)})}{\det(A)}=\frac{20}{20}=1.
\]
Por tanto, $\boxed{(x,y,z)=(1,1,1)}$. Al sustituir en las ecuaciones,
\[
    1+2(1)+1=4,\qquad 2(1)-1+1=2,\qquad 3(1)+1-2(1)=2.
\]

\section*{4. Determinante de matrices semejantes (2 puntos)}
\noindent\textbf{Rúbrica (2 puntos).}
\begin{itemize}[itemsep=2pt]
    \item Aplicar el determinante a $A=PBP^{-1}$ y usar su propiedad
          multiplicativa: 0,75 puntos.
    \item Justificar $\det(P^{-1})=1/\det(P)$ y concluir
          $\det(A)=\det(B)$: 0,75 puntos.
    \item Obtener $\det(B)=5$ e identificar la opción (b): 0,5 puntos.
\end{itemize}

De $A=PBP^{-1}$ y de la propiedad multiplicativa del determinante,
\[
    \det(A)=\det(P)\det(B)\det(P^{-1}).
\]
Como $P$ es invertible, $\det(P)\ne0$ y
$\det(P^{-1})=1/\det(P)$. En consecuencia,
\[
    \det(A)=\det(B),\qquad \boxed{\det(B)=5.}
\]
La respuesta correcta es la \textbf{opción (b)}.

\end{document}
